\documentclass[10pt,letterpaper]{article}
\usepackage{cornell-notes}

\title{Clause 31.05.02.02.06: Class Ios Base Init}
\author{}
\date{}
\setDocTitle{Clause 31.05.02.02.06: Class Ios Base Init}
\setDocSubtitle{C++ 2024}
\setDocOwner{C++ 2024 Cornell Notes}
\setDocDate{\today}
\setCornellCollection{C++ 2024 Cornell Notes}
\setCornellUnitType{Clause}
\setCornellUnitNumber{31.05.02.02.06}
\setCornellUnitTitle{Class Ios Base Init}
\hypersetup{pdftitle={Clause 31.05.02.02.06: Class Ios Base Init},pdfsubject={Cornell notes},pdfauthor={}}

\begin{document}
\maketitle
\begin{CornellOverviewBox}{Overview}
\textbf{Classification:} Standard library - Normative\\[2pt]
\textbf{Learning objective:} Understand the rules, terminology, and semantic consequences associated with Class ios\_base::Init.
\end{CornellOverviewBox}\begin{CornellSummaryBox}{Summary}
{\sffamily\bfseries\color{Accent} Summary}\par\vspace{3pt}Section 31.5.2.2.6 focuses on Class ios\_base::Init. Apply the specified interface together with its mandates, effects, result conditions, exception behavior, and complexity constraints. When applying it, preserve the distinction between mandatory requirements, recommendations, permissions, and behavior deliberately left undefined, unspecified, or implementation-defined.
\end{CornellSummaryBox}

\end{document}
